% !TEX root = ../../ctfp-print.tex

\lettrine[lhang=0.17]{古}{代ギリシャの劇作家エウリピデスはかつて言った：}「人はつきあ
う仲間に似るものだ。」私たちは関係性によって定義される。これは圏論においてこれ以上な
いほど真実だ。圏の中で特定の対象を際立たせたいなら、他の対象（および自分自身）との関
係のパターンを記述することによってのみ可能だ。これらの関係は射によって定義される。

圏論には、対象をその関係の観点から定義するための\newterm{普遍的構成}と呼ばれる一般的
な構成がある。これを行う一つの方法は、対象と射で構成された特定の形状であるパターンを
選び、圏内のそのすべての出現を探すことだ。それが十分に一般的なパターンであり、圏が大
きければ、たくさんのヒットが得られる可能性が高い。コツはそれらのヒットの間に何らかの
ランキングを確立し、最良の適合と考えられるものを選ぶことだ。

このプロセスはウェブ検索の方法を思わせる。クエリはパターンのようなものだ。非常に一般
的なクエリは大きな\emph{再現率}をもたらす：多くのヒット。一部は関連しているかもしれ
ず、他は関連していないかもしれない。無関連なヒットを排除するためにクエリを絞り込む。
それにより\emph{精度}が上がる。最終的に、検索エンジンはヒットをランク付けし、望まし
くは、関心のある結果がリストの上位に来る。

\section{始対象}

最も単純な形状は単一の対象だ。明らかに、この形状のインスタンスは与えられた圏の対象と
同じ数だけある。選ぶものがたくさんある。何らかのランキングを確立し、この階層の頂点に
立つ対象を見つけようとする必要がある。使えるのは射だけだ。射を矢印と考えるなら、圏の
一端から他端への矢印の全体的な流れがあるかもしれない。これは順序付き圏、例えば半順序
において当てはまる。対象$a$から対象$b$への矢印（射）があれば「対象$a$の方がより始的
だ」と言うことで、対象の先行関係というその概念を一般化できる。それから、すべての他の
対象への矢印を持つ対象として\emph{始}対象を定義したい。明らかにそのような対象が存在
する保証はないが、それでも構わない。より大きな問題はそのような対象が多すぎる可能性が
あることだ：再現率は良いが精度が欠けている。解決策は順序付き圏からヒントを得ることだ
――任意の2つの対象間の矢印は高々1つ：別の対象より小さいか等しいという方法はただ一つ
しかない。これが始対象の定義につながる：

\begin{quote}
  \textbf{始対象}は圏内の任意の対象へ向かう射がただ一つだけある対象だ。
\end{quote}

\begin{figure}[H]
  \centering
  \includegraphics[width=0.4\textwidth]{images/initial.jpg}
\end{figure}

\noindent
しかし、それでも始対象の一意性は保証されない（存在する場合）。しかし次に良いものは保
証される：\newterm{同型を除いた}一意性だ。同型は圏論において非常に重要なので、間もな
く話す。今のところ、同型を除いた一意性が始対象の定義における「ある」（the）の使用を
正当化することに同意しよう。

いくつか例を示す。半順序集合（しばしば\newterm{ポセット}と呼ばれる）の始対象はその最
小元だ。一部のポセットは始対象を持たない――例えば、射として以下または等しい関係を持つ
すべての整数（正と負）の集合のように。

集合と関数の圏では、始対象は空集合だ。空集合はHaskellの型\code{Void}に対応すること
（C++には対応する型がない）と、\code{Void}から他の任意の型への一意な多相関数は
\code{absurd}と呼ばれることを覚えているだろう：

\src{snippet01}
\code{Void}を型の圏における始対象にするのは、この射のファミリーだ。

\section{終対象}

単一対象のパターンを続けよう。しかし対象のランク付け方法を変えよう。対象$b$から対象
$a$への射があれば（方向の逆転に注意）「対象$a$の方がより終的だ」と言おう。圏内の他の
任意の対象よりもより終的な対象を探す。再び、一意性を主張する：

\begin{quote}
  \textbf{終対象}は圏内の任意の対象からそこへ来る射がただ一つだけある対象だ。
\end{quote}

\begin{figure}[H]
  \centering
  \includegraphics[width=0.4\textwidth]{images/final.jpg}
\end{figure}

\noindent
そして再び、終対象は同型を除いて一意だ。これはすぐに示す。しかしまずいくつか例を見て
みよう。ポセットでは、終対象（存在すれば）は最大元だ。集合の圏では、終対象は単集合だ。
単集合についてはすでに話した――C++の\code{void}型とHaskellのユニット型\code{()}に対応
する。それは値がただ一つしかない型だ――C++では暗黙的、Haskellでは\code{()}で明示的に
示される。また、任意の型からユニット型への純粋関数はただ一つしかないことも確立した：

\src{snippet02}
だから終対象のすべての条件が満たされる。

この例では一意性の条件が重要だ。なぜなら、すべての集合（実際は空集合以外のすべて）か
ら入ってくる射を持つ他の集合があるからだ。例えば、すべての型に対して定義されたブール
値関数（述語）がある：

\src{snippet03}
しかし\code{Bool}は終対象ではない。\code{Bool}値関数はすべての型からさらに少なくとも
もう一つある（\code{Void}を除く、そこでは両方の関数が\code{absurd}に等しい）：

\src{snippet04}
一意性を主張することで、終対象の定義をただ一つの型に絞り込むのにちょうど良い精度が
得られる。

\section{双対性}

始対象と終対象を定義した方法の対称性に気づかずにはいられない。2つの間の唯一の違いは
射の方向だった。任意の圏$\cat{C}$に対して、すべての矢印を逆にすることで\newterm{双対
圏}$\cat{C}^\mathit{op}$を定義できることがわかる。双対圏は、合成を同時に再定義する限
り、圏のすべての要件を自動的に満たす。元の射$f \Colon a \to b$と$g \Colon b \to c$が
$h = g \circ f$で$h \Colon a \to c$に合成されるなら、逆にした射$f^\mathit{op} \Colon
b \to a$と$g^\mathit{op} \Colon c \to b$は$h^\mathit{op} = f^\mathit{op} \circ
g^\mathit{op}$で$h^\mathit{op} \Colon c \to a$に合成される。そして恒等射を逆にすること
は（ダジャレ注意！）何もしない（no-op）。

双対性は圏論において非常に重要な性質だ。なぜなら、圏論を研究するすべての数学者の生産
性を2倍にするからだ。考案した構成のすべてに対して、その双対がある；証明したすべての定
理に対して、一つ無料で得られる。双対圏での構成はしばしば「余（co）」が前置される。だ
から積と余積、モナドと余モナド、錐と余錐、極限と余極限などがある。ただし余余モナドは
ない。矢印を2回逆にすると元の状態に戻るからだ。

したがって、終対象は双対圏における始対象だ。

\section{同型}

プログラマーとして、等しさを定義することが非自明なタスクだということはよく知っている。
2つのオブジェクトが等しいとはどういう意味か？メモリ上の同じ場所を占有しなければなら
ないか（ポインタ等式）？それともすべての成分の値が等しければ十分か？複素数は、一方が
実部と虚部で表現され、他方が係数と角度で表現されているとき等しいか？数学者は等しさの
意味を解明しているだろうと思うかもしれないが、解明していない。彼らも等しさの複数の競
合する定義という同じ問題を抱えている。命題的等式、内包的等式、外延的等式、ホモトピー
型理論でのパスとしての等式がある。そして同型というより弱い概念、さらに弱い同値がある。

直感は、同型な対象は同じように見える――同じ形を持つということだ。一方の対象のすべての
部分が、一対一のマッピングで他方の対象のある部分に対応することを意味する。私たちの観
測手段が判断できる限り、2つの対象は互いの完全なコピーだ。数学的には、対象$a$から対象
$b$へのマッピングがあり、対象$b$から対象$a$へのマッピングがあり、それらが互いの逆で
あることを意味する。圏論ではマッピングを射に置き換える。同型は逆を持つ射；つまり、一
方が他方の逆である射のペアだ。

合成と恒等の観点から逆を理解する：射$g$が射$f$の逆であるのは、それらの合成が恒等射で
ある場合だ。2つの射を合成する方法が2つあるので、実際には2つの等式がある：

\src{snippet05}
始（終）対象が同型を除いて一意だと言ったとき、任意の2つの始（終）対象は同型であるこ
とを意味した。これは実際に見て取るのが簡単だ。2つの始対象$i_{1}$と$i_{2}$があるとし
よう。$i_{1}$が始的なので、$i_{1}$から$i_{2}$への一意な射$f$がある。同様に、$i_{2}$
が始的なので、$i_{2}$から$i_{1}$への一意な射$g$がある。これらの2つの射の合成は何か？

\begin{figure}[H]
  \centering
  \includegraphics[width=0.4\textwidth]{images/uniqueness.jpg}
  \caption{この図のすべての射は一意だ。}
\end{figure}

\noindent
合成$g \circ f$は$i_{1}$から$i_{1}$への射でなければならない。しかし$i_{1}$は始的なの
で、$i_{1}$から$i_{1}$への射は一つしかない。圏の中にいるので、$i_{1}$から$i_{1}$への
恒等射があることを知っており、場所は一つしかないので、それでなければならない。したがっ
て$g \circ f$は恒等に等しい。同様に、$i_{2}$から$i_{2}$への射は一つしかないので、
$f \circ g$も恒等に等しくなければならない。これにより$f$と$g$は互いの逆でなければなら
ないことが証明される。したがって任意の2つの始対象は同型だ。

この証明では始対象からそれ自身への射の一意性を使ったことに注意しよう。それなしでは
「同型を除いた」部分を証明できなかった。しかしなぜ$f$と$g$の一意性が必要なのか？なぜ
なら始対象は同型を除いて一意なだけでなく、\emph{一意な}同型を除いて一意だからだ。原
則として、2つの対象の間に複数の同型があるかもしれないが、ここではそうではない。この
「一意な同型を除いた一意性」がすべての普遍的構成の重要な性質だ。

\section{積}

次の普遍的構成は積の構成だ。2つの集合の直積とは何かを知っている：ペアの集合だ。しか
し積集合とその構成要素の集合を結びつけるパターンは何か？それを解明できれば、他の圏に
一般化できる。

言えるのは、積からそれぞれの構成要素への2つの関数、射影があるということだ。Haskellで
はこれらの2つの関数は\code{fst}と\code{snd}と呼ばれ、それぞれペアの第1成分と第2成分
を取り出す：

\src{snippet06}

\src{snippet07}
ここで関数は引数のパターンマッチングで定義されている：任意のペアにマッチするパターン
は\code{(x, y)}で、その成分を変数\code{x}と\code{y}に展開する。

これらの定義はワイルドカードを使うことでさらに簡略化できる：

\src{snippet08}
C++ではテンプレート関数を使う：

\begin{snip}{cpp}
template<typename A, typename B>
A fst(pair<A, B> const & p) {
    return p.first;
}
\end{snip}
この一見非常に限られた知識を持って、集合の圏における対象と射のパターンを定義し、2つの
集合$a$と$b$の積の構成につながるものを探してみよう。このパターンは対象$c$と、それを
$a$と$b$にそれぞれ結ぶ2つの射$p$と$q$から成る：

\src{snippet09}

\begin{figure}[H]
  \centering
  \includegraphics[width=0.3\textwidth]{images/productpattern.jpg}
\end{figure}

\noindent
このパターンに合うすべての$c$が積の候補として考えられる。たくさんあるかもしれない。

\begin{figure}[H]
  \centering
  \includegraphics[width=0.4\textwidth]{images/productcandidates.jpg}
\end{figure}

\noindent
例えば、構成要素として2つのHaskellの型\code{Int}と\code{Bool}を選び、それらの積の候
補のサンプルを取ってみよう。

一つ目：\code{Int}。\code{Int}は\code{Int}と\code{Bool}の積の候補と考えられるか？はい、
できる――そしてその射影はこうだ：

\src{snippet10}
かなり貧弱だが、基準を満たしている。

もう一つ：\code{(Int, Int, Bool)}。3つの要素のタプル、つまり三つ組だ。これを正当な候
補にする2つの射がある（三つ組のパターンマッチングを使う）：

\src{snippet11}
最初の候補は小さすぎる――積の\code{Int}次元しかカバーしていない；2番目は大きすぎる――
\code{Int}次元が余分に複製されている、ということに気づいたかもしれない。

しかし普遍的構成の別の部分：ランキングをまだ探っていない。パターンの2つのインスタンス
を比較したい。一方の候補対象$c$とその2つの射影$p$と$q$を、別の候補対象$c'$とその2つ
の射影$p'$と$q'$と比較したい。$c'$から$c$への射$m$があれば$c$は$c'$より「良い」と言
いたいが、それは弱すぎる。また、その射影も$c'$の射影より「良い」または「より普遍的で
ある」ことを望む。$m$を使って$p$と$q$から$p'$と$q'$を再構成できることを意味する：

\src{snippet12}

\begin{figure}[H]
  \centering
  \includegraphics[width=0.4\textwidth]{images/productranking.jpg}
\end{figure}

\noindent
これらの等式を見る別の方法は、$m$が$p'$と$q'$を\emph{因子分解}するということだ。これ
らの等式が自然数にあり、ドットが掛け算だと仮定してみよう：$m$は$p'$と$q'$が共有する
公因数だ。

直感を養うために、2つの標準的な射影\code{fst}と\code{snd}を持つペア\code{(Int, Bool)}
が確かに前に提示した2つの候補より\emph{良い}ことを示そう。

\begin{figure}[H]
  \centering
  \includegraphics[width=0.4\textwidth]{images/not-a-product.jpg}
\end{figure}

\noindent
最初の候補の\code{m}はこうだ：

\src{snippet13}
確かに、2つの射影\code{p}と\code{q}は次のように再構成できる：

\src{snippet14}
2番目の例の\code{m}も同様に一意に決まる：

\src{snippet15}
\code{(Int, Bool)}が2つの候補のどちらよりも良いことを示せた。逆が真でない理由を見て
みよう。\code{p}と\code{q}から\code{fst}と\code{snd}を再構成するのに役立つ\code{m'}
を見つけられるか？

\src{snippet16}
最初の例では、\code{q}は常に\code{True}を返し、第2成分が\code{False}であるペアがあ
ることを知っている。\code{q}から\code{snd}を再構成できない。

2番目の例は異なる：\code{p}または\code{q}のどちらを実行した後でも十分な情報を保持し
ているが、\code{fst}と\code{snd}を因子分解する方法が複数ある。\code{p}も\code{q}も
三つ組の第2成分を無視するので、\code{m'}はそこに何でも入れることができる：

\src{snippet17}

または
\src{snippet18}
などなど。

まとめると、2つの射影\code{p}と\code{q}を持つ任意の型\code{c}に対して、それらを因子
分解する\code{c}から直積\code{(a, b)}への一意な\code{m}がある。実際、それは単に
\code{p}と\code{q}をペアに組み合わせる：

\src{snippet19}
これにより直積\code{(a, b)}が最良の適合となる。つまりこの普遍的構成は集合の圏で機能
する。任意の2つの集合の積を選び出す。

今、集合について忘れ、同じ普遍的構成を使って任意の圏での2つの対象の積を定義しよう。
そのような積は常に存在するわけではないが、存在する場合、一意な同型を除いて一意だ。

\begin{quote}
  2つの対象$a$と$b$の\textbf{積}は、2つの射影を持つ対象$c$であり、2つの射影を持つ他
  の任意の対象$c'$に対して、それらの射影を因子分解する$c'$から$c$への一意な射$m$が
  存在するものだ。
\end{quote}

\noindent
因子分解関数\code{m}を2つの候補から生成する（高階）関数は\newterm{因子分解子}と呼ば
れることがある。この場合、それは次の関数だ：

\src{snippet20}

\section{余積}

圏論のすべての構成と同様に、積には双対があり、余積と呼ばれる。積のパターンで矢印を逆
にすると、2つの\emph{入射}\code{i}と\code{j}（$a$と$b$から$c$への射）を持つ対象$c$に
なる。

\src{snippet21}

\begin{figure}[H]
  \centering
  \includegraphics[width=0.4\textwidth]{images/coproductpattern.jpg}
\end{figure}

\noindent
ランキングも逆転する：$c$から$c'$への射$m$があって、入射を因子分解するなら（入射$i'$
と$j'$を持つ対象$c'$に対して）、対象$c$は対象$c'$より「良い」：

\src{snippet22}

\begin{figure}[H]
  \centering
  \includegraphics[width=0.4\textwidth]{images/coproductranking.jpg}
\end{figure}

\noindent
他の任意のパターンへそこから連なる一意な射を持つような「最良の」対象が余積と呼ばれ、
存在する場合、一意な同型を除いて一意だ。

\begin{quote}
  2つの対象$a$と$b$の\textbf{余積}は、2つの入射を持つ対象$c$であり、2つの入射を持つ
  他の任意の対象$c'$に対して、それらの入射を因子分解する$c$から$c'$への一意な射$m$が
  存在するものだ。
\end{quote}

\noindent
集合の圏では、余積は2つの集合の\emph{非交和}だ。$a$と$b$の非交和の要素は$a$の要素か
$b$の要素だ。2つの集合が重なっていれば、非交和には共通部分の2つのコピーが含まれる。
非交和の要素は出自を指定するタグが付いていると考えることができる。

プログラマーにとって、余積を型の観点から理解する方が簡単だ：それは2つの型のタグ付き
共用体だ。C++は共用体をサポートするが、タグ付きではない。プログラムで共用体のどのメン
バーが有効かを何らかの方法で追跡しなければならないことを意味する。タグ付き共用体を作
るにはタグ――列挙型――を定義し、共用体と組み合わせる必要がある。例えば、\code{int}と
\code{char const *}のタグ付き共用体はこのように実装できる：

\begin{snip}{cpp}
struct Contact {
    enum { isPhone, isEmail } tag;
    union { int phoneNum; char const * emailAddr; };
};
\end{snip}
2つの入射はコンストラクタまたは関数として実装できる。例えば、関数\code{PhoneNum}とし
ての最初の入射はこうだ：

\begin{snip}{cpp}
Contact PhoneNum(int n) {
    Contact c;
    c.tag = isPhone;
    c.phoneNum = n;
    return c;
}
\end{snip}
整数を\code{Contact}に入射する。

タグ付き共用体は\newterm{バリアント}とも呼ばれ、C++17標準に非常に一般的なバリアントの
実装\code{std::variant}がある。

Haskellでは、縦棒でデータコンストラクタを区切ることで任意のデータ型をタグ付き共用体
に組み合わせることができる。\code{Contact}の例はこの宣言に変換される：

\src{snippet23}
ここで、\code{PhoneNum}と\code{EmailAddr}はコンストラクタ（入射）としても、パターン
マッチングのタグとしても機能する（これについては後ほど詳しく説明する）。例えば、電話
番号を使ってコンタクトを構成する方法はこうだ：

\src{snippet24}
Haskellに組み込まれたプリミティブなペアとして実装された積の標準的な実装とは異なり、
余積の標準的な実装は\code{Either}と呼ばれるデータ型であり、標準Preludeで次のように定
義されている：

\src{snippet25}
2つの型\code{a}と\code{b}でパラメータ化されており、2つのコンストラクタを持つ：型
\code{a}の値を取る\code{Left}と、型\code{b}の値を取る\code{Right}だ。

積の因子分解子を定義したように、余積のものも定義できる。候補型\code{c}と2つの候補入
射\code{i}と\code{j}が与えられたとき、\code{Either}の因子分解子は因子分解関数を生成
する：

\src{snippet26}

\section{非対称性}

2組の双対的な定義を見た：終対象の定義は始対象の定義から矢印の方向を逆にすることで得
られる；同様に、余積の定義は積の定義から得られる。それでも集合の圏では始対象は終対象
と大きく異なり、余積は積と大きく異なる。後で積は乗算のように振る舞い、終対象が1の役
割を果たすことを見る；一方余積は和のように振る舞い、始対象がゼロの役割を果たす。特に
有限集合では、積のサイズは個々の集合のサイズの積であり、余積のサイズはサイズの和だ。

これは集合の圏が矢印の反転に関して対称的でないことを示している。

空集合は任意の集合への一意な射（\code{absurd}関数）を持つが、戻ってくる射はないことに
注意しよう。単集合はすべての集合から来る一意な射を持つが、すべての集合（空集合を除く）
への出て行く射も\emph{また}持つ。前に見たように、終対象からのこれらの出て行く射は他の
集合の要素を選ぶ非常に重要な役割を果たす（空集合には要素がないので、選ぶものがない）。

単集合と積の関係こそが余積と区別するものだ。ユニット型\code{()}で表される単集合を、積
パターンのさらに別の――格段に劣る――候補として使うことを考えよう。2つの射影\code{p}と
\code{q}：単集合から各構成要素の集合への関数を持たせよう。それぞれの集合から具体的な
要素を選ぶ。積が普遍的なので、候補である単集合から積への（一意な）射\code{m}もある。
この射は積集合から要素を選ぶ――具体的なペアを選ぶ。また2つの射影を因子分解する：

\src{snippet27}
単集合の唯一の要素であるシングルトン値\code{()}に作用するとき、これらの2つの等式はこ
うなる：

\src{snippet28}
\code{m ()}は\code{m}が選ぶ積の要素なので、これらの等式は、\code{p}が最初の集合から
選ぶ要素\code{p ()}が\code{m}が選ぶペアの第1成分であることを教えてくれる。同様に、
\code{q ()}は第2成分に等しい。これは積の要素が構成要素の集合からの要素のペアであると
いう私たちの理解と完全に一致している。

余積にはそのような単純な解釈がない。余積から要素を抽出しようとして、候補として単集合
を試みることができるが、そこでは2つの射影が出て行くのではなく、2つの入射が入ってくる。
それらは源について何も教えてくれない（実際、入力パラメータを無視することを見た）。余積
から単集合への一意な射もそうだ。集合の圏は始対象の方向から見たときと終端から見たとき
では非常に異なって見える。

これは集合の本質的な性質ではなく、$\Set$における射として使っている関数の性質だ。関数
は一般的に非対称だ。説明しよう。

関数はその定義域集合のすべての要素に対して定義されなければならない（プログラミングでは
\newterm{全}関数と呼ぶ）が、終域全体をカバーする必要はない。いくつかの極端な例を見た：
単集合からの関数――終域でただ一つの要素だけを選ぶ関数。（実際は空集合からの関数が本当
の極端だ。）定義域のサイズが終域のサイズよりずっと小さい場合、そのような関数を定義域
を終域に埋め込むものとしてよく考える。例えば、単集合からの関数をその単一の要素を終域
に埋め込むものとして考えることができる。これらを\newterm{埋め込み}関数と呼ぶが、数学
者は反対のものに名前をつけることを好む：終域をぎっしり埋める関数は\newterm{全射的}ま
たは\newterm{onto}と呼ばれる。

非対称性のもう一つの源は、関数が定義域集合の多くの要素を終域の一つの要素にマップする
ことが許されていることだ。それらは圧縮できる。極端な場合は集合全体を単集合にマップす
る関数だ。まさにそれを行う多相的な\code{unit}関数を見た。圧縮は合成によって複合され
るだけだ。2つの圧縮関数の合成は個々の関数よりさらに圧縮される。数学者は圧縮しない関
数に名前をつけている：\newterm{単射的}または\newterm{一対一}と呼ぶ。

もちろん、埋め込みでも圧縮でもない関数もある。それらは\newterm{全単射}と呼ばれ、可逆
なので本当に対称的だ。集合の圏では、同型は全単射と同じだ。

\section{課題}

\begin{enumerate}
  \tightlist
  \item
        終対象が一意な同型を除いて一意であることを示せ。
  \item
        ポセットにおける2つの対象の積は何か？ヒント：普遍的構成を使え。
  \item
        ポセットにおける2つの対象の余積は何か？
  \item
        Haskellの\code{Either}に相当するものを好きな言語（Haskell以外）でジェネリック
        型として実装せよ。
  \item
        \code{Either}が2つの入射を持つ\code{int}より「良い」余積であることを示せ：

        \begin{snip}{cpp}
int i(int n) { return n; }
int j(bool b) { return b ? 0: 1; }
\end{snip}

        ヒント：関数

        \begin{snip}{cpp}
int m(Either const & e);
\end{snip}

        を定義して\code{i}と\code{j}を因子分解せよ。
  \item
        前の問題を続けて：2つの入射\code{i}と\code{j}を持つ\code{int}が\code{Either}
        より「良い」ことができない理由をどのように論証するか？
  \item
        引き続き：これらの入射についてはどうか？

        \begin{snip}{cpp}
int i(int n) {
    if (n < 0) return n;
    return n + 2;
}

int j(bool b) { return b ? 0: 1; }
\end{snip}
  \item
        \code{int}と\code{bool}の余積の劣った候補を考案せよ。それは\code{Either}への
        複数の許容可能な射を許すため、\code{Either}より良くなれない。
\end{enumerate}

\section{参考文献}

\begin{enumerate}
  \tightlist
  \item
        The Catsters,
        \urlref{https://www.youtube.com/watch?v=upCSDIO9pjc}{Products and
          Coproducts}（積と余積）動画。
\end{enumerate}
